\chapter{Objectives Specification and Work Environment}
\chaplabel{ch:specification}

\section*{Introduction}
Chapter \ref{ch:context} identified a gap between the debug and trace capability a microcontroller implements and the capability a developer can reach through available tools. This chapter states what the project undertook to do about that gap, the requirements derived from it, and the hardware and software used to produce the result.

\section{Project Objectives and Deliverables}
\label{sec:objectives}
The internship produced the three deliverables of Table \ref{tab:deliverables}. Chapter \ref{ch:implementation} presents the chronology.

\begin{table}[H]
    \centering
    \caption{Deliverables}
    \label{tab:deliverables}
    \begin{tabularx}{\linewidth}{| @{}>{\hsize=0.5\hsize\bfseries}Y | >{\hsize=1.5\hsize}Y@{} |}
        \hline
        Deliverable & \textbf{Description} \\
        \hline
        Application note & The instruction trace infrastructure of STM32 microcontrollers, a reference implementation for decoding the trace such an infrastructure produces, and guidance on working with instruction trace features on vendor IDEs. \\
        \hline
        Proof of Concept & A demonstration that instruction trace is obtainable on an STM32 device using an on-chip trace sink. \\
        \hline
        Debugger & A debugger with instruction trace features. \\
        \hline
    \end{tabularx}
\end{table}

\subsection{Scope}
\label{subsec:scope}

\begin{table}[H]
    \centering
    \caption{Scope of the three deliverables}
    \label{tab:scope}
    \begin{tabularx}{\linewidth}{| @{}>{\hsize=0.6\hsize\bfseries}Y | >{\hsize=1.4\hsize}Y@{} |}
        \hline
        Deliverable & \textbf{Scope} \\
        \hline
        \multirow{3}{*}{Application note} & Instruction trace infrastructure of every STM32 family that implements it \\
        \cline{2-2}
        & Keil \textmu Vision with the ULINKpro probe \\
        \cline{2-2}
        & IAR Embedded Workbench with the I-jet Trace probe \\
        \hline
        \multirow{3}{*}{Proof of Concept} & ETMv4 implementations \\
        \cline{2-2}
        & On-chip trace sink \\
        \cline{2-2}
        & Instruction trace \\
        \hline
        \multirow{4}{*}{Debugger} & STM32 MCUs \\
        \cline{2-2}
        & Breakpoints and variable inspection \\
        \cline{2-2}
        & Peripheral inspection, memory views and disassembly \\
        \cline{2-2}
        & Instruction trace \\
        \hline
    \end{tabularx}
\end{table}

\section{Requirements Analysis}
\label{sec:requirements}

\subsection{Functional Requirements}
\label{subsec:fr}

\begin{table}[H]
    \centering
    \caption{Functional requirements}
    \label{tab:fr}
    \begin{tabularx}{\linewidth}{| @{}>{\hsize=0.4\hsize\bfseries}Y | >{\hsize=1.6\hsize}Y@{} |}
        \hline
        ID & \textbf{Requirement} \\
        \hline
        FR1 & Configure the trace unit and the on-chip trace sink of the target through the debug connection. \\
        \hline
        FR2 & Start and stop trace capture in step with the run state of the program, without a separate command from the developer. \\
        \hline
        FR3 & Extract the captured trace over the same debug connection that carries every other debug operation. \\
        \hline
        FR4 & Decode the extracted byte stream into trace elements. \\
        \hline
        FR5 & Reconstruct the sequence of executed instructions from those elements together with the program image. \\
        \hline
        FR6 & Attribute each reconstructed instruction to a function, a source file and a line. \\
        \hline
        FR7 & Represent a discontinuity in the capture explicitly, rather than joining the two sides of it. \\
        \hline
        FR8 & Present the reconstruction inside a development tool during a live debug session. \\
        \hline
        FR9 & Provide the run control primitives of Section \ref{sec:primitives}, namely halting, resuming, stepping, breakpoints and watchpoints. \\
        \hline
        FR10 & Read and write memory and processor registers. \\
        \hline
        FR11 & Present the registers of a peripheral by name and by field, using a description of the device. \\
        \hline
        FR12 & Provide disassembly and source-level inspection of variables. \\
        \hline
        FR13 & Reset the target and restore the debug session to a state consistent with the reset device. \\
        \hline
    \end{tabularx}
\end{table}

\subsection{Non-Functional Requirements}
\label{subsec:nfr}

\begin{table}[H]
    \centering
    \caption{Non-functional requirements}
    \label{tab:nfr}
    \begin{tabularx}{\linewidth}{| @{}>{\hsize=0.4\hsize\bfseries}Y | >{\hsize=1.6\hsize}Y@{} |}
        \hline
        ID & \textbf{Requirement} \\
        \hline
        NFR1 & Require no hardware beyond the debug probe already present on the development board, and no trace pins. \\
        \hline
        NFR2 & Operate on Windows and on Linux. \\
        \hline
        NFR3 & Integrate with existing development tools. \\
        \hline
    \end{tabularx}
\end{table}

\subsection{Traceability to Deliverables}
\label{subsec:traceability}

\begin{table}[H]
    \centering
    \caption{Traceability of requirements to deliverables}
    \label{tab:traceability}
    \begin{tabularx}{\linewidth}{| @{}>{\hsize=1.2\hsize\bfseries}Y | >{\hsize=0.9\hsize}Y | >{\hsize=0.9\hsize}Y@{} |}
        \hline
        Deliverable & \textbf{Functional} & \textbf{Non-functional} \\
        \hline
        Application note & FR1 to FR5 & NFR1 \\
        \hline
        Proof of Concept & FR1 to FR8 & NFR1 \\
        \hline
        Debugger & FR8 to FR13 & NFR2, NFR3 \\
        \hline
    \end{tabularx}
\end{table}

\section{Hardware Environment}
\label{sec:hardware}

\subsection{Microcontroller: STM32H7S3L8}
\label{subsec:target}
The main device used during the development period of the project is an STM32H7S3L8, a member of the STM32H7RS family. It carries an ARM Cortex-M7 processor \cite{cortexm7_trm} running at up to 600 MHz with a memory protection unit, a double-precision floating-point unit and separate instruction and data caches of 32 kilobytes each. The device provides 64 kilobytes of internal flash memory and 620 kilobytes of internal SRAM, and is supplied in a TFBGA225 package \cite{ds_h7s3}.

Figure \ref{fig:target_debug} shows the debug and trace components of the device and the buses that connect them.

\begin{figure}[H]
    \centering
    \includegraphics[width=\linewidth]{stm32h7rs_debug_infrastructure.pdf}
    \caption[Debug and trace infrastructure of the target]{Debug and trace infrastructure of the target, adapted from \cite{rm0477}}
    \label{fig:target_debug}
\end{figure}

\subsection{Development Board: NUCLEO-H7S3L8}
\label{subsec:board}
The microcontroller is used on a NUCLEO-H7S3L8. It exposes a MIPI20 connector for external debug tools, which carries the serial wire and JTAG signals together with a four-bit trace port \cite{um3276}.

\subsection{Debug Probe: STLINK-V3EC}
\label{subsec:probe}
The probe is the STLINK-V3EC embedded on the board, which supports both SWD and JTAG \cite{um3276}. It connects to the host through a USB Type-C connector and additionally provides a virtual COM port (VCP) and mass storage programming.

\subsection{Trace Resources Available on the Target}
\label{subsec:trace_resources}

\begin{table}[H]
    \centering
    \caption{Trace components of the target device}
    \label{tab:trace_resources}
    \begin{tabularx}{\linewidth}{| @{}>{\hsize=0.55\hsize\bfseries}Y | >{\hsize=0.8\hsize}Y | >{\hsize=1.65\hsize}Y@{} |}
        \hline
        Name & \textbf{Base address} & \textbf{Note} \\
        \hline
        ETM & 0xE0041000 & Implementation of the Cortex-M7 ETM, configured for instruction trace only \cite{etm_m7_trm}. \\
        \hline
        TMC & 0xE00F4000 & ETF configuration, providing two kilobytes of memory. \\
        \hline
        TPIU & 0xE00F5000 & - \\
        \hline
        Funnel & 0xE00F3000 & Funnels the ETM trace data into the TMC. It also carries the stream of the instrumentation macrocell \cite{rm0477}, which this work does not use. \\
        \hline
    \end{tabularx}
\end{table}

\section{Software Environment}
\label{sec:software}

\subsection{Development Environment}
\label{subsec:devenv}

\begin{table}[H]
    \centering
    \caption{Development environment}
    \label{tab:devenv}
    \begin{tabularx}{\linewidth}{| @{}>{\hsize=0.85\hsize\bfseries}Y | >{\hsize=1.15\hsize}Y@{} |}
        \hline
        Element & \textbf{Choice} \\
        \hline
        Host platforms & Linux and Windows \\
        \hline
        Backend language & C++23 \\
        \hline
        Frontend language & TypeScript \\
        \hline
        Compiler & Clang \\
        \hline
        Build system & CMake, Ninja generator \\
        \hline
        Frontend bundler & esbuild \\
        \hline
        Editor & Visual Studio Code \\
        \hline
    \end{tabularx}
\end{table}

\subsection{External Software Components}
\label{subsec:extsw}

\begin{table}[H]
    \centering
    \caption{External software components}
    \label{tab:software}
    \begin{tabularx}{\linewidth}{| @{}>{\hsize=0.65\hsize\bfseries}Y | >{\hsize=0.4\hsize}Y | >{\hsize=1.45\hsize}Y | >{\hsize=1.5\hsize}Y@{} |}
        \hline
        Component & \textbf{Version} & \textbf{Description} & \textbf{Justification} \\
        \hline
        OpenOCD & Fork of 0.12.0 & Open-source on-chip debugger \cite{openocd}. & Supplies the groundwork the trace drivers extend. \\
        \hline
        LLVM & 22 & Compiler infrastructure and debugger, of which the machine code layer \cite{llvm_mc} and the debugger \cite{lldb} are used. & Supplies the core debug functionality and the analysis of the program image. \\
        \hline
        OpenCSD & 1.8.3 & Library for decoding CoreSight trace \cite{opencsd}. & Simplifies the trace decoding process. \\
        \hline
        Xerces-C & 3.2.4 & XML parser. & Reads the CMSIS-SVD device descriptions. \\
        \hline
        CodeSynthesis XSD & 4.2.0 & Generator of C++ bindings from an XML schema. & Avoids writing the CMSIS-SVD bindings by hand. \\
        \hline
        Visual Studio Code & 1.125 & Editor hosting the frontend extension. & Already implements the client side of the Debug Adapter Protocol. \\
        \hline
    \end{tabularx}
\end{table}

\section{Conclusion}
This chapter has stated the three deliverables of the project, the requirements derived from them, and the environment in which they were produced. Two facts from it govern everything that follows.

The first is the hardware. The target implements both routes by which trace leaves a processor, and the bench can use only one of them, since the probe available cannot capture a trace port. The on-chip sink is therefore the route the work takes, and its capacity of two kilobytes is the constraint against which the results are read.

The second is the requirement set. The functional requirements describe a pipeline running from the configuration of a trace unit to the display of an executed source line, and the non-functional requirements constrain it to hardware already present and to an existing development tool. Chapter \ref{ch:architecture} describes the structure built to satisfy them.
